\section{From primitive entry kernels to operator majorants}
\label{sec:operator-certificates}

\subsection{Tail-free entrywise Abel identity}

Set
\[
 M(x)=\sum_{r\le x}\mu(r)
\]
and, for a zero-extended finite kernel \(F\),
\[
 F^{[d]}(x,y)=F(dx,dy),
 \qquad
 \Delta_{12}F(x,y)
 =
 F(x,y)-F(x+1,y)-F(x,y+1)+F(x+1,y+1).
\]
Define the location-sensitive weighted plaquette cost
\begin{equation}
 \cC_\square(F)
 =
 \sum_{x,y\ge1}
 |M(x)M(y)|\,|\Delta_{12}F(x,y)|.
 \label{eq:entry-plaquette-cost}
\end{equation}
For comparison with the divisor lift of TPC--73, also define its exact
large-diagonal tail
\begin{equation}
 \cT_D(F)
 =
 \sum_{d>D}|\mu(d)|
 \sum_{x,y\ge1}|\mu(dx)\mu(dy)F(dx,dy)|.
 \label{eq:entry-diagonal-tail}
\end{equation}

\begin{theorem}[Primitive entry Abel certificate]
\label{thm:primitive-entry-Abel}
For every \(m\ne n\),
\begin{align}
 \widetilde H_{mn}
 &=
 \sum_{x,y\ge1}
 M(x)M(y)\Delta_{12}L_{mn}^{0}(x,y),
 \label{eq:full-entry-Abel}\\
 \widetilde E_{mn}
 &=
 \sum_{x,y\ge1}
 M(x)M(y)\Delta_{12}L_{mn}^{E}(x,y).
 \label{eq:erasure-entry-Abel}
\end{align}
Consequently,
\begin{align}
 |\widetilde H_{mn}|
 &\le \cC_\square(L_{mn}^{0}),
 \label{eq:full-entry-majorant}\\
 |\widetilde E_{mn}|
 &\le \cC_\square(L_{mn}^{E}),
 \label{eq:erasure-entry-majorant}\\
 |\widetilde R_{mn}|
 &=|L_{mn}(b_{mn},a_{mn})|.
 \label{eq:residual-entry-exact-magnitude}
\end{align}
Moreover, for \(d\ge2\),
\[
 (L_{mn}^{0})^{[d]}\equiv0,\qquad
 (L_{mn}^{E})^{[d]}\equiv0.
\]
Thus the auxiliary diagonal tails of the support-faithful entry kernels
vanish for every cutoff \(D\ge1\).
\end{theorem}

\begin{proof}
Both kernels are finitely supported on the primitive lattice.  Apply
two-dimensional summation by parts to
\eqref{eq:full-direct-plane} and
\eqref{eq:erasure-direct-plane}, then take absolute values.  If \(d\ge2\),
the point \((dx,dy)\) has greatest common divisor at least \(d\), so both
primitive-supported kernels vanish there.  The residual magnitude follows
from \eqref{eq:residual-direct-point}.
\end{proof}

\begin{remark}[No cancellation theorem yet]
Equations \eqref{eq:full-entry-Abel} and
\eqref{eq:erasure-entry-Abel} are exact.  They do not assert that either
weighted plaquette cost is smaller than the literal entry mass.  The
prime number theorem statement \(M(x)=o(x)\) does not by itself supply a
uniform fixed-power saving over the growing row-pair family.
\end{remark}

\subsection{The nonnegative entry matrix}

For \(m<n\), define
\begin{align}
 P^H_{mn}
 &=\cC_\square(L_{mn}^{0}),
 \label{eq:full-Perron-majorant}\\
 P^E_{mn}
 &=\cC_\square(L_{mn}^{E}),
 \label{eq:erasure-Perron-majorant}\\
 P^R_{mn}
 &=|L_{mn}(b_{mn},a_{mn})|.
 \label{eq:residual-Perron-majorant}
\end{align}
Set
\[
 P^\chi_{nm}=P^\chi_{mn},\qquad
 P^\chi_{mm}=0
 \quad(\chi\in\{H,E,R\}).
\]
This upper-triangle definition makes all three matrices symmetric without
an orientation convention.  It agrees with the reversed kernels by
\cref{thm:reversed-pair-conjugacy}, since transposition and conjugation
preserve \(\cC_\square\), while
\((b_{nm},a_{nm})=(a_{mn},b_{mn})\).

\begin{theorem}[Literal Perron operator certificates]
\label{thm:literal-Perron-certificates}
The exact physical matrices satisfy
\begin{align}
 \boxed{\|H\|=\|\widetilde H\|\le\|P^H\|=\rho(P^H),}
 \label{eq:full-Perron-certificate}\\
 \boxed{\|E\|=\|\widetilde E\|\le\|P^E\|=\rho(P^E),}
 \label{eq:erasure-Perron-certificate}\\
 \boxed{\|R\|=\|\widetilde R\|\le\|P^R\|=\rho(P^R).}
 \label{eq:residual-Perron-certificate}
\end{align}
For every \(c=(c_m)\in(0,\infty)^{\cM_M}\),
\begin{equation}
 \rho(P^H)
 \le
 \max_m\frac{\sum_nP^H_{mn}c_n}{c_m},
 \qquad
 \rho(P^H)
 =
 \inf_{c_m>0}
 \max_m\frac{\sum_nP^H_{mn}c_n}{c_m}.
 \label{eq:weighted-Collatz}
\end{equation}
The analogous statements hold for \(P^E\) and \(P^R\).
\end{theorem}

\begin{proof}
The entry inequalities in
\cref{thm:primitive-entry-Abel} give
\(|\widetilde H_{mn}|\le P^H_{mn}\), and similarly for \(E,R\).
For arbitrary vectors \(x,y\),
\[
 |\langle x,\widetilde H y\rangle|
 \le
 \langle |x|,P^H|y|\rangle
 \le
 \|P^H\|\,\|x\|_2\|y\|_2.
\]
Hence \(\|\widetilde H\|\le\|P^H\|\).  The majorant is symmetric and
nonnegative, so its operator norm is its Perron root.  The weighted bound
and the reducible-safe infimum formula are the Collatz--Wielandt
characterization \citep{BermanPlemmons1994}.  The other matrices are
identical.
\end{proof}

\begin{corollary}[Schur and row-star forms]
\label{cor:Schur-row-star}
Let
\[
 \Delta_H
 =
 \max_m\#\{n:P^H_{mn}>0\}.
\]
Then
\begin{equation}
 \|H\|
 \le
 \max_m\sum_nP^H_{mn}
 \le
 \sqrt{\Delta_H}\,
 \max_m\left(\sum_n(P^H_{mn})^2\right)^{1/2}.
 \label{eq:row-star-majorant}
\end{equation}
The weighted form \eqref{eq:weighted-Collatz} may be strictly sharper than
the unweighted maximum row sum.
\end{corollary}

\begin{corollary}[A conditional normalized-Gram consequence]
\label{cor:conditional-Gram}
For a family indexed by \(X\), if a literal certificate proves
\[
 \rho(P_X^H)\le\eta_X<1,
\]
then
\[
 1-\eta_X
 \le
 \lambda_{\min}(G_{{\rm nat},X})
 \le
 \lambda_{\max}(G_{{\rm nat},X})
 \le
 1+\eta_X.
\]
This implication is deterministic.  The hypothesis is not proved here.
\end{corollary}

\begin{proof}
Use \(G_{\rm nat}=I+H\) and
\eqref{eq:full-Perron-certificate}.
\end{proof}

\subsection{Completion freedom returns at the estimate level}

The physical entry samples only primitive sites.  Therefore the analytic
completion gauge of TPC--73 can be applied separately to every
\(L_{mn}^{0}\).  A nonprimitive completion may reduce its \(d=1\)
variation, but all higher diagonal slices and the exact tail of that same
completion must be restored.  To obtain an operator theorem, the complete
family of completion rules must be declared before the later coefficient
extremizer is chosen, must respect
\eqref{eq:kernel-Hermitian-reversal}, and must produce one controlled
majorant matrix.  Pairwise optimization with no uniform summability is
not an L2 result.
